%%%%%%%%%%%%%%%%%%%%%%%%%%%%%%%%%%%%%%%%%%%%%%%%%%%%%%%%%%%%%%%%%%%%%%
% Lớp 10 | Phương pháp tọa độ trong mặt phẳng (𝑂𝑥𝑦) | Ba đường conic trong mặt phẳng tọa độ
% Dạng 0H9?5-8: Phương trình chính tắc của parabol
% Số câu: 2  (sinh tự động bởi tach_cau.py; từ nay thầy sửa trực tiếp file này)
%%%%%%%%%%%%%%%%%%%%%%%%%%%%%%%%%%%%%%%%%%%%%%%%%%%%%%%%%%%%%%%%%%%%%%

%%%%%%%%%%%%%%%% ĐÚNG/SAI %%%%%%%%%%%%%%%%

% nguồn: MaHoaNH_DungSai_Lop10
\begin{ex}%[0H9H5-8]%[ma:e31ae04f]
	Trong các mệnh đề dưới  đây, mệnh đề nào đúng, mệnh đề nào sai?
	 \choiceTF
		 {Parabol có tiêu điểm $F(8;0)$ có phương trình là $y^2=16x$}
		{\True Parabol có khoảng cách từ tiêu điểm  đến đường chuẩn bằng $4$ có phương trình là $y^2=8x$}
		 {Tọa độ điểm $M\in (P)\colon y^2=4x$ và $MF=3$ là $M(2;8)$}
		{Tọa độ điểm  $M\in (P)\colon y^2=2x$ và $MF=3$ là $M(2;6)$}
	\loigiai{ 
		\begin{itemchoice}
			\itemch Vì tiêu điểm $F(8;0)$ nên $\dfrac{p}{2}=8\Leftrightarrow p=16$. Phương trình chính tắc  của $(P)\colon y^2=32x$.
			\itemch Vì khoảng cách từ tiêu điểm đến đường  chuẩn bằng $4$ nên $p=4$. Phương trình chính tắc của $(P)$ là $y^2=8x$.
			\itemch Gọi $M(x_0;y_0)\in (P)$ nên ta  có
			$$\heva{&y_0^2=4x_0\\&|x_0+2|=3}\Rightarrow  x_0=1\Rightarrow y_0=\pm 2.$$
			 Vậy $M(2;8)\notin (P)\colon y^2=4x$.
			\itemch Gọi  $M(x_0;y_0)\in (P)$ nên ta có
			$$\heva{&y_0^2=2x_0\\&|x_0+1|=3}\Rightarrow  x_0=2\Rightarrow y_0=\pm 2.$$
			Vậy $M(2;6)\notin (P)\colon y^2=2x$. 
		\end{itemchoice}
	} 
\end{ex}


%%%%%%%%%%%%%%%% TRẢ LỜI NGẮN %%%%%%%%%%%%%%%%

% nguồn: MaHoaNH_TraLoiNgan_Lop10
\begin{bt}%[0H9H5-8]%[ma:5226ef47]
	Cho parabol $(P)\colon y^2=2px$ ($p>0$). Biết $(P)$ có đường chuẩn $\Delta$ song song và cách đường thẳng $d\colon x=2$ một khoảng bằng $5$. Khi đó giá trị  của $p$ bằng bao nhiêu?
	\shortans{6} 
	 \loigiai{
		Vì $(P)$ có phương  trình chính tắc là $y^2=2px$ ($p>0$) nên phương trình đường chuẩn $\Delta$ có dạng $x=-\dfrac{p}{2}$. Theo giả thiết
		$$\mathrm{d}(d,\Delta)=5\Rightarrow  \left|\dfrac{-p}{2}-2 \right|=5\Leftrightarrow\hoac{& \dfrac{-p}{2}-2=5\\&\dfrac{-p}{2}-2=-5 } \Leftrightarrow \hoac{&p=-14<0 \text{ (loại)}\\& p=6>0.}$$
		Vậy  $p=6$.
	 }
\end{bt}
